
\subsubsection{2.1 Iterative Code Repair}

Self-Debug \citep{chen2023selfdebug} introduced execution-feedback loops for LLM code generation, using test results and stack traces to guide iterative refinement. Self-Refine \citep{madaan2023selfrefine} generalized this to broader refinement tasks using LLM self-feedback. These methods establish execution feedback as the baseline for repair loops.

LDB \citep{zhong2024ldb} extended this paradigm with block-by-block verification, improving localization of errors. However, these approaches rely on execution traces---they identify where code fails but not why.

\subsubsection{2.2 Static Analysis for LLMs}

Blyth et al. (2025) represent the closest prior work. They integrated static analysis feedback (pylint, bandit) into LLM pipelines and demonstrated improvements in code quality metrics: security issues reduced from 40\% to 13\%, and readability problems reduced from 80\% to 11\%. Notably, they did not measure functional correctness (pass@k). Their contribution establishes that LLMs can interpret static warnings, but whether this improves benchmark performance remains unknown.

Chen et al. (2024) proposed combining testing and static analysis feedback for code generation, though their evaluation focused on quality metrics rather than pass@k.

\subsubsection{2.3 The Gap Addressed}

\begin{table}[htbp]
\centering
\resizebox{\columnwidth}{!}{%
\begin{tabular}{llll}
\toprule
Prior Work & Feedback Type & Metric & Pass@k Measured \\
\midrule
Self-Debug (Chen et al., 2023) & Execution & Functional correctness & Yes \\
Self-Refine (Madaan et al., 2023) & LLM self-feedback & Task completion & Yes \\
Blyth et al. (2025) & Static analysis & Quality (security, readability) & No \\
**This work** & Static + Execution & Methodology baseline & Pending \\
\bottomrule
\end{tabular}
}
\end{table}

No published study measures static analysis impact on pass@k. This work provides the methodology and baseline to enable such evaluation.

\subsubsection{2.4 Benchmarks}

HumanEval \citep{chen2021humaneval} contains 164 Python programming problems and is widely used for code generation evaluation. MBPP \citep{austin2021mbpp} contains 974 problems. Pass@k metrics measure functional correctness---the fraction of problems solved within k attempts. This study focuses on HumanEval, with MBPP extension as future work.
